\chapter{Cloud Deployment, DevOps, dan Observability}
\label{ch:deployment}

\section*{Tujuan Pembelajaran}
\begin{itemize}
  \item Mengemas satu layanan sebagai \textit{container} yang setelannya
        datang dari \textit{environment}, lalu menjalankan dua versi dari
        artefak yang sama persis.
  \item Menjalankan rilis \textit{blue-green}, \textit{canary}, dan
        \textit{rollback} lewat satu \textit{reverse proxy}, lalu mengukur
        pembagian permintaan yang benar-benar terjadi.
  \item Menginstrumentasi layanan dengan \textit{log}, metrik, dan
        \textit{trace}, lalu menilainya terhadap SLO sebagai alarm yang
        keluar dengan status gagal.
\end{itemize}

\section{Pendahuluan}

Delapan bab sebelumnya membangun dan memperbaiki aplikasi di komputer
sendiri. Bab ini menanyakan apa yang berubah ketika aplikasi itu dijalankan
di tempat lain, dilayani untuk orang lain, dan diperbarui tanpa boleh mati.

Tiga pertanyaan praktis muncul di sana. Bagaimana satu artefak yang sama
dijalankan di lingkungan yang berbeda tanpa diubah isinya. Bagaimana versi
baru dinaikkan tanpa membuat seluruh pengunjung menanggung risikonya.
Bagaimana keadaan sistem diketahui dari luar, tanpa menebak dari keluhan
pengguna.

Demonya sengaja kecil: dua versi layanan yang sama di belakang satu
\textit{reverse proxy}, seluruhnya dengan Python pustaka standar, ditambah
satu \textit{image} container berisi satu file. Referensi utamanya adalah
aturan \textit{twelve-factor app} \cite{twelve-factor}, W3C Trace Context
\cite{w3c-trace-context} untuk menyatukan \textit{trace}, dan uraian SLO
beserta \textit{error budget} dari buku SRE Google \cite{beyer2016sre}.

\section{Container dan Twelve-Factor App}

\textit{Container} mengemas aplikasi beserta seluruh yang dibutuhkannya
menjadi satu artefak yang dapat dijalankan di mana saja. Yang membuatnya
berguna bukan pengemasannya saja, melainkan janji bahwa artefak yang diuji
dan artefak yang dijalankan di produksi benar-benar sama.

\textit{Twelve-factor app} \cite{twelve-factor} adalah daftar aturan yang
membuat janji itu dapat ditepati. Tabel \ref{tab:12factor} memuat lima aturan
yang paling terasa pada demo bab ini.

\begin{table}[!htb]
  \centering
  \caption{Lima aturan \textit{twelve-factor} dan wujudnya pada demo Bab 9}
  \label{tab:12factor}
  \begin{tabularx}{\textwidth}{lX}
    \hline
    \textbf{Aturan} & \textbf{Wujudnya di sini} \\
    \hline
    Setelan dari \textit{environment} & \texttt{VERSI}, \texttt{PORT}, dan \texttt{HOST} dibaca saat proses mulai \\
    \hline
    Satu artefak, banyak lingkungan & Dua versi berjalan dari \textit{image} yang sama, berbeda hanya di \textit{environment} \\
    \hline
    Proses tanpa ingatan tetap & Metrik hanya hidup selama proses, dan hilang saat proses diganti \\
    \hline
    \textit{Port binding} & Layanan mengikat portnya sendiri, tanpa server web di luarnya \\
    \hline
    \textit{Log} sebagai aliran peristiwa & Tiap permintaan satu baris JSON ke stdout, tanpa menulis file \\
    \hline
  \end{tabularx}
\end{table}

Aturan pertama itu tampak sepele sampai dilanggar. Pada demo ini, layanan
semula mengikat ke \texttt{127.0.0.1}, dan itu bekerja dengan baik di
komputer sendiri. Begitu dimasukkan ke \textit{container}, port yang
dipublikasikan tidak pernah terjawab, karena di dalam \textit{container}
alamat itu berarti \textit{container} itu sendiri. Perbaikannya satu baris:
alamat ikat ikut dibaca dari \textit{environment}, dan \textit{image}
menyetelnya menjadi \texttt{0.0.0.0}.

Tabel \ref{tab:container} memuat angka \textit{image} dan container tersebut.

\begin{table}[!htb]
  \centering
  \caption{Ukuran dan waktu \textit{image} layanan Bab 9, hasil pengukuran
           nyata}
  \label{tab:container}
  \begin{tabularx}{\textwidth}{Xr}
    \hline
    \textbf{Ukuran} & \textbf{Hasil} \\
    \hline
    Ukuran \textit{image} & 119 MB \\
    \hline
    Lama \textit{build} pertama, termasuk mengunduh \textit{base image} & 129,7 s \\
    \hline
    Lama \textit{build} ulang, seluruh \textit{layer} tersimpan & 9,5 s \\
    \hline
    Lama sampai siap melayani sesudah \texttt{docker run} & 0,5 s \\
    \hline
  \end{tabularx}
\end{table}

Selisih 129,7 detik dengan 9,5 detik itu bukan keajaiban, melainkan akibat
urutan di dalam \texttt{Dockerfile}. \textit{Layer} yang jarang berubah
ditulis lebih dahulu, dan \textit{layer} yang sering berubah ditulis paling
akhir, sehingga \textit{build} berikutnya cukup mengulang bagian terakhirnya
\cite{docker-build-best}.

\section{Konfigurasi dan Manajemen Secret}

Setelan dan \textit{secret} sama-sama datang dari luar kode, tetapi
diperlakukan berbeda. Bab 7 sudah menetapkan aturannya: setelan biasa boleh
masuk file yang ikut repositori, sedangkan \textit{secret} hanya lewat
\textit{environment} atau \textit{secret manager}.

Bab ini menambahkan satu pembedaan lagi, yaitu setelan yang boleh berubah
tanpa menyalakan ulang. Tabel \ref{tab:setelan9} memuat ketiganya beserta
contohnya pada demo bab ini.

\begin{table}[!htb]
  \centering
  \caption{Tiga jenis nilai dari luar kode dan cara memperlakukannya}
  \label{tab:setelan9}
  \begin{tabularx}{\textwidth}{lXl}
    \hline
    \textbf{Jenis} & \textbf{Contoh di Bab 9} & \textbf{Berubahnya} \\
    \hline
    Setelan proses & \texttt{VERSI}, \texttt{PORT}, \texttt{HOST} & Saat proses mulai \\
    \hline
    Setelan operasional & Bobot rilis di \berkas{rilis.json} & Kapan saja, tanpa menyalakan ulang \\
    \hline
    \textit{Secret} & Kunci sesi, kata sandi basis data & Lewat \textit{environment} saja \\
    \hline
  \end{tabularx}
\end{table}

Baris tengah yang membuat rilis pada bab ini mungkin. Bobot rilis dibaca
proxy pada tiap permintaan, sehingga mengubah satu angka sudah cukup untuk
memindahkan lalu lintas, tanpa satu proses pun dimatikan.

\section{CI/CD dan Infrastructure as Code}

\textit{Continuous Integration} sudah dipakai sejak Bab 5 sebagai gerbang
yang menolak perubahan cacat. \textit{Continuous Delivery} melanjutkannya:
perubahan yang lolos gerbang dikemas dan disiapkan untuk dirilis, tanpa
langkah manual yang diingat-ingat.

Urutan tahapnya mengikuti aturan yang sama dengan Bab 6, yaitu yang paling
murah dan paling sering gagal dijalankan lebih dahulu. Listing
\ref{lst:ci9} memuat potongan contoh \textit{pipeline} pada bab ini.

\begin{lstlisting}[language=yaml, caption={Potongan contoh pipeline, tersedia di \berkas{sources/chapter09/ci-contoh.yml}}, label={lst:ci9}]
      - name: Analisis statis
        run: ruff check sources/chapter09

      - name: Uji unit dan integrasi
        run: pytest -q sources/chapter06

      - name: Bangun image
        run: docker build -t awe9-layanan sources/chapter09
\end{lstlisting}

\textit{Infrastructure as Code} (IaC) membawa gagasan yang sama ke
infrastrukturnya: keadaan yang diinginkan ditulis sebagai file, bukan
dikerjakan lewat panel. File \berkas{sources/chapter09/docker-compose.yml}
adalah bentuk paling sederhananya. Di dalamnya tertulis dua layanan yang
dibangun dari \textit{image} yang sama dan hanya berbeda \textit{environment},
dan siapa pun yang menjalankannya mendapat susunan yang sama persis.

\section{Reverse Proxy serta Deployment ke Cloud dan Edge}

\textit{Reverse proxy} berdiri di depan satu atau beberapa layanan, menerima
seluruh permintaan, lalu meneruskannya. Karena seluruh lalu lintas
melewatinya, di sanalah tempat paling murah untuk memutuskan siapa yang
melayani apa. Gambar \ref{fig:topologi9} memperlihatkan susunan pada demo
bab ini.

\begin{figure}[!htb]
  \centering
  \begin{tikzpicture}[
    font=\scriptsize,
    kotak/.style={draw, rounded corners, minimum height=0.9cm,
                  minimum width=2.2cm, align=center, font=\scriptsize\bfseries},
    panah/.style={-{Stealth[length=2mm]}, thick}
  ]
    \node[kotak, fill=blue!8] (klien) at (0,0) {Klien};
    \node[kotak, fill=orange!12] (proxy) at (3.4,0) {Reverse proxy\\{\tiny\mdseries rilis.json}};
    \node[kotak, fill=blue!8] (biru) at (7.4,0.9) {Versi biru\\{\tiny\mdseries port 8061}};
    \node[kotak, fill=green!10] (hijau) at (7.4,-0.9) {Versi hijau\\{\tiny\mdseries port 8062}};

    \draw[panah] (klien) -- (proxy);
    \draw[panah] (proxy.east) -- ++(0.4,0) |- (biru.west);
    \draw[panah] (proxy.east) -- ++(0.4,0) |- (hijau.west);
    \node[font=\tiny, gray, anchor=east] at (6.9,0.72) {90 persen};
    \node[font=\tiny, gray, anchor=east] at (6.9,-0.72) {10 persen};
    \node[font=\tiny, gray, anchor=north, align=center] at (3.4,-1.3)
      {membuat trace id,\\mencatat tujuannya};
  \end{tikzpicture}
  \caption{Seluruh permintaan melewati proxy, sehingga pembagian rilis dan
           penandaan \textit{trace} cukup dilakukan di satu tempat}
  \label{fig:topologi9}
\end{figure}

Di \textit{cloud}, peran proxy itu biasanya dipegang \textit{load balancer}
milik penyedia, dan di \textit{edge} dipegang jaringan titik terdekat seperti
dibahas Bab 8. Bentuknya berbeda, tetapi keputusannya sama: memilih tujuan,
menambahkan penanda, dan mencatat hasilnya.

\section{Rilis Blue-Green, Canary, dan Rollback}

Tiga nama ini menjawab satu pertanyaan: berapa banyak pengunjung yang
menanggung risiko versi baru. \textit{Blue-green} memindahkan seluruhnya
sekaligus. \textit{Canary} memindahkan sebagian kecil lebih dahulu.
\textit{Rollback} mengembalikan seluruhnya ketika angkanya memburuk.

Listing \ref{lst:pilih} memperlihatkan seluruh mekanismenya, yaitu empat
baris yang memilih tujuan menurut bobot.

\begin{lstlisting}[style=PythonStyle, caption={Pembagian rilis dalam empat baris, tersedia di \berkas{sources/chapter09/proxy.py}}, label={lst:pilih}]
def pilih_tujuan(rilis, nomor_permintaan):
  """Memilih versi tujuan menurut bobot pada file rilis.

  Pemilihannya berurutan, bukan acak, supaya hasilnya dapat diulang: dari
  tiap seratus permintaan, tepat sebanyak bobot yang menuju versi kedua.
  """
  bobot_kedua = rilis["bobot_persen_kedua"]
  if nomor_permintaan % 100 < bobot_kedua:
    return rilis["versi_kedua"]
  return rilis["versi_pertama"]
\end{lstlisting}

Versi hijau pada demo ini sengaja dibuat cacat: pekerjaannya empat kali lebih
lama, dan satu dari dua puluh lima permintaannya dijawab 500. Tabel
\ref{tab:rilis} memuat ketiga tahap rilisnya, diukur dari header
\texttt{X-Dilayani} pada tiap \textit{response}.

\begin{table}[!htb]
  \centering
  \caption{Tiga tahap rilis pada demo Bab 9, hasil pengukuran nyata}
  \label{tab:rilis}
  \begin{tabularx}{\textwidth}{Xrrrr}
    \hline
    \textbf{Tahap} & \textbf{Minta} & \textbf{Biru} & \textbf{Hijau} & \textbf{Gagal} \\
    \hline
    \textit{Canary} 10 persen & 500 & 450 & 50 & 2 \\
    \hline
    \textit{Blue-green}, seluruhnya hijau & 100 & 0 & 100 & 4 \\
    \hline
    \textit{Rollback}, kembali ke biru & 100 & 100 & 0 & 0 \\
    \hline
  \end{tabularx}
\end{table}

Tiga baris itu merangkum seluruh alasan \textit{canary} ada. Pada tahap
pertama, cacat versi hijau hanya mengenai 50 dari 500 permintaan, dan hanya
dua yang gagal. Pada tahap kedua, cacat yang sama mengenai seluruhnya, dan
empat dari seratus permintaan gagal. Pada tahap ketiga, satu angka di file
rilis dikembalikan dan kegagalannya hilang seketika, tanpa satu proses pun
dimatikan.

\section{OpenTelemetry: Log, Metrik, dan Trace}

Tiga keluaran instrumentasi menjawab tiga pertanyaan yang berbeda
\cite{otel-docs}. \textit{Log} menjawab apa yang terjadi pada satu peristiwa.
Metrik menjawab bagaimana keadaan seluruhnya dari waktu ke waktu.
\textit{Trace} menjawab ke mana saja satu permintaan pergi.

Tabel \ref{tab:tiga-pilar} memuat wujud ketiganya pada demo bab ini.

\begin{table}[!htb]
  \centering
  \caption{Tiga keluaran instrumentasi dan wujudnya pada demo Bab 9}
  \label{tab:tiga-pilar}
  \begin{tabularx}{\textwidth}{lXl}
    \hline
    \textbf{Keluaran} & \textbf{Wujudnya} & \textbf{Menjawab} \\
    \hline
    \textit{Log} & Satu baris JSON per permintaan ke stdout & Apa yang terjadi \\
    \hline
    Metrik & Teks format Prometheus di \texttt{/metrics} & Bagaimana keadaannya \\
    \hline
    \textit{Trace} & \textit{Trace id} di header \texttt{traceparent} & Ke mana perginya \\
    \hline
  \end{tabularx}
\end{table}

Metriknya memakai format teks Prometheus \cite{prometheus-exposition}, dan
sebaran waktunya disimpan sebagai histogram. Cara itu hemat: berapa pun
jumlah permintaannya, yang disimpan hanya hitungan per ember. Harganya,
persentil yang dihitung darinya selalu jatuh pada salah satu batas ember,
bukan pada angka persis.

Papan operasional pada bab ini membaca \texttt{/metrics} kedua versi, lalu
menampilkannya berdampingan. Tabel \ref{tab:papan} memuat keluarannya sesudah
tahap \textit{canary}.

\begin{table}[!htb]
  \centering
  \caption{Papan operasional sesudah 500 permintaan \textit{canary}, hasil
           pengukuran nyata}
  \label{tab:papan}
  \begin{tabularx}{\textwidth}{Xrrrrr}
    \hline
    \textbf{Versi} & \textbf{Minta} & \textbf{Gagal} & \textbf{\% gagal} & \textbf{Rata ms} & \textbf{p95 ms} \\
    \hline
    biru & 450 & 0 & 0,0 & 5,1 & 10 \\
    \hline
    hijau & 50 & 2 & 4,0 & 20,1 & 25 \\
    \hline
  \end{tabularx}
\end{table}

Pertanyaan berikutnya, bagaimana satu permintaan ditelusuri ketika melewati
dua proses. Jawabannya \textit{trace id}, yaitu penanda yang dibuat sekali di
proxy lalu diteruskan lewat header \texttt{traceparent} menurut W3C Trace
Context \cite{w3c-trace-context}. Kedua pihak mencatat penanda yang sama,
sehingga baris lognya dapat disatukan:

\begin{lstlisting}[language=bash, caption={Menyatukan log dua proses dengan satu trace id, memakai \berkas{sources/chapter09/proxy.py}}, label={lst:trace9}]
grep -h "300a032b60413b55e53009cc0660f240" proxy.log hijau.log

# Keluarannya, dirapikan:
#   hijau    /api/katalog   status 500   20.2 ms
#   proxy    /api/katalog   status 500   22.2 ms
\end{lstlisting}

Dua baris itu menjawab pertanyaan yang tidak dapat dijawab metrik maupun
\textit{log} sendirian: kegagalan tersebut berasal dari layanan hijau, bukan
dari proxy, dan selisih 2 milidetik di antara keduanya adalah biaya
penerusan.

\section{Dasar SLI dan SLO}

\textit{Service Level Indicator} (SLI) adalah angka yang mewakili pengalaman
pengguna, misalnya persentase permintaan yang berhasil.
\textit{Service Level Objective} (SLO) adalah batas yang disepakati untuk
angka itu, misalnya 99 persen \cite{beyer2016sre}.

Selisih antara 100 persen dan SLO disebut \textit{error budget}. Dengan SLO
99 persen, dari 500 permintaan ada jatah 5 permintaan yang boleh gagal.
Jatah itulah yang membuat rilis mungkin: selama budgetnya belum habis, versi
baru boleh dinaikkan.

Tabel \ref{tab:slo} memuat penilaian alarm pada bab ini sesudah tahap
\textit{canary}.

\begin{table}[!htb]
  \centering
  \caption{Penilaian SLO sesudah tahap \textit{canary}, hasil pengukuran
           nyata}
  \label{tab:slo}
  \begin{tabularx}{\textwidth}{Xrrl}
    \hline
    \textbf{Tujuan} & \textbf{SLO} & \textbf{Terukur} & \textbf{Verdict} \\
    \hline
    biru: ketersediaan persen & 99,0 & 100,0 & lolos \\
    \hline
    biru: p95 milidetik & 50 & 10 & lolos \\
    \hline
    hijau: ketersediaan persen & 99,0 & 96,0 & LANGGAR \\
    \hline
    hijau: p95 milidetik & 50 & 25 & lolos \\
    \hline
  \end{tabularx}
\end{table}

Alarm itu keluar dengan status 1, sama seperti gerbang pada Bab 5 sampai Bab
7, sehingga dapat dipasang di penjadwal maupun di \textit{pipeline} sesudah
rilis. Yang menarik justru baris p95 versi hijau: angkanya 25 milidetik, dua
setengah kali lipat versi biru, tetapi masih di bawah batas 50 milidetik.
Pelan belum tentu melanggar, dan itulah gunanya menuliskan batasnya lebih
dahulu.

Satu keterbatasan demo ini perlu disebut. Metriknya dihitung sejak proses
mulai, sedangkan SLO yang sungguhan dihitung atas jendela bergerak, misalnya
dua puluh delapan hari terakhir. Akibatnya, pada demo ini kegagalan lama
tidak pernah kedaluwarsa, dan satu-satunya cara mengosongkannya adalah
menyalakan ulang prosesnya.

\section{Ringkasan}

Satu artefak dijalankan di banyak lingkungan ketika seluruh setelannya datang
dari luar kode. Demo bab ini menjalankan dua versi dari \textit{image} yang
sama persis, berbeda hanya pada \textit{environment}-nya, dan satu
pelanggaran kecil terhadap aturan itu langsung terasa: layanan yang mengikat
ke \texttt{127.0.0.1} bekerja di komputer sendiri, lalu tidak pernah terjawab
begitu masuk \textit{container}.

Rilis adalah keputusan tentang berapa banyak pengunjung yang menanggung
risiko. Pada tahap \textit{canary} sepuluh persen, cacat versi hijau hanya
mengenai 50 dari 500 permintaan dengan dua kegagalan. Ketika seluruh lalu
lintas dipindahkan, cacat yang sama menghasilkan empat kegagalan dari seratus
permintaan. \textit{Rollback} mengembalikannya dengan mengubah satu angka,
tanpa mematikan satu proses pun.

Keadaan sistem diketahui dari tiga keluaran yang menjawab pertanyaan berbeda.
\textit{Log} menjawab apa yang terjadi, metrik menjawab bagaimana keadaan
seluruhnya, dan \textit{trace} menjawab ke mana permintaan itu pergi. Pada
demo ini ketiganya cukup ditulis dengan pustaka standar, dan \textit{trace
id} dari header \texttt{traceparent} sudah cukup menyatukan dua proses.

Angka instrumentasi baru berguna ketika ada batas yang disepakati. SLO 99
persen memberi jatah kegagalan yang jelas, dan alarm pada bab ini menilai
kedua versi terhadapnya: versi biru lolos dengan 100 persen, versi hijau
melanggar dengan 96 persen, lalu prosesnya keluar dengan status 1.

Pelan belum tentu melanggar, dan itulah alasan batasnya ditulis lebih dahulu.
Versi hijau dua setengah kali lebih lambat daripada versi biru, tetapi p95
25 milidetiknya masih di bawah batas 50 milidetik. Tanpa batas yang
disepakati, angka seperti itu hanya menjadi bahan perdebatan.

\section{Latihan}

Seluruh file yang dipakai pada latihan ini tersedia di
\berkas{sources/chapter09/}, dihitung relatif terhadap akar repositori.
Latihan membutuhkan Python 3, dan Docker hanya pada latihan terakhir.

Seluruh pengukuran diarahkan ke layanan yang berjalan di komputer sendiri,
dengan beban setara satu pengguna yang mengirim permintaan berurutan. Hasil
tiap mahasiswa akan berbeda pada bagian waktu, dan yang dilaporkan adalah
pola perbandingannya beserta penjelasannya.

Penyiapannya dilakukan sekali, seperti pada Listing \ref{lst:c9-siapkan}.

\begin{lstlisting}[language=bash, caption={Menyalakan dua versi layanan beserta proxynya, memakai \berkas{sources/chapter09/aplikasi.py}}, label={lst:c9-siapkan}]
cd sources/chapter09
VERSI=biru PORT=8061 python3 aplikasi.py > biru.log &
VERSI=hijau PORT=8062 python3 aplikasi.py > hijau.log &
python3 proxy.py > proxy.log &
\end{lstlisting}

\subsection*{Latihan 1: Menjalankan Rilis Canary}

Latihan ini membuktikan bahwa bobot rilis menentukan berapa banyak pengunjung
yang menanggung risiko versi baru.

\begin{enumerate}
  \item Pastikan \texttt{bobot\_persen\_kedua} di
        \berkas{sources/chapter09/rilis.json} bernilai 10, lalu kirim 200
        permintaan seperti pada Listing \ref{lst:c9-lat1-jalan}.
  \item Catat pembagian yang terjadi beserta jumlah kegagalannya ke Tabel
        \ref{tab:c9-lat1-hasil}. Tabel \ref{tab:c9-lat1-contoh}
        memperlihatkan contoh yang sudah terisi.
  \item Ubah bobotnya menjadi 50, kirim 200 permintaan lagi, lalu catat
        kembali. Perhatikan bahwa proxy tidak perlu dinyalakan ulang.
  \item Jawab tiga pertanyaan berikut. Pertama, berapa permintaan yang
        menanggung risiko pada tiap bobot? Kedua, mengapa pembagiannya tepat
        dan tidak meleset seperti pembagian acak? Ketiga, apa yang menjadi
        pertimbangan ketika memilih bobot pertama sebuah rilis?
\end{enumerate}

\begin{lstlisting}[language=bash, caption={Mengukur pembagian rilis, tersedia di \berkas{sources/chapter09/ukur-rilis.py}}, label={lst:c9-lat1-jalan}]
python3 ukur-rilis.py 200

# Keluarannya:
# Permintaan   : 200
#   biru     :  180 permintaan (90.0%)
#   hijau    :   20 permintaan (10.0%)
#   status 200 :  200 permintaan
# p50          : 6.7 ms
# p95          : 21.9 ms
\end{lstlisting}

\begin{table}[!htb]
  \centering
  \caption{Contoh hasil Latihan 1 yang sudah terisi}
  \label{tab:c9-lat1-contoh}
  \begin{tabularx}{\textwidth}{Xrrrr}
    \hline
    \textbf{Bobot hijau} & \textbf{Minta} & \textbf{Biru} & \textbf{Hijau} & \textbf{Gagal} \\
    \hline
    10 persen & 500 & 450 & 50 & 2 \\
    \hline
    100 persen & 100 & 0 & 100 & 4 \\
    \hline
  \end{tabularx}
\end{table}

\begin{table}[!htb]
  \centering
  \caption{Lembar isian hasil pengukuran Latihan 1}
  \label{tab:c9-lat1-hasil}
  \begin{tabularx}{\textwidth}{Xrrrr}
    \hline
    \textbf{Bobot hijau} & \textbf{Minta} & \textbf{Biru} & \textbf{Hijau} & \textbf{Gagal} \\
    \hline
    10 persen & \dots & \dots & \dots & \dots \\
    \hline
    50 persen & \dots & \dots & \dots & \dots \\
    \hline
    100 persen & \dots & \dots & \dots & \dots \\
    \hline
    0 persen, sesudah \textit{rollback} & \dots & \dots & \dots & \dots \\
    \hline
  \end{tabularx}
\end{table}

\subsection*{Latihan 2: Membaca Papan Operasional}

Latihan ini membuktikan bahwa perbedaan dua versi terbaca dari metrik, tanpa
perlu membaca kodenya.

\begin{enumerate}
  \item Lihat teks metrik mentah salah satu versi, lalu kenali ketiga jenis
        metriknya, seperti pada Listing \ref{lst:c9-lat2-jalan}.
  \item Jalankan papan operasionalnya, lalu catat angka kedua versi ke Tabel
        \ref{tab:c9-lat2-hasil}. Tabel \ref{tab:c9-lat2-contoh}
        memperlihatkan contoh yang sudah terisi.
  \item Kirim 200 permintaan lagi, lalu jalankan papannya sekali lagi dan
        perhatikan angka mana yang bertambah dan angka mana yang tetap.
  \item Jawab dua pertanyaan berikut. Pertama, mengapa p95 pada papan selalu
        jatuh pada angka yang sama, misalnya 10 atau 25? Kedua, apa yang
        hilang bila sebaran waktu disimpan sebagai rata-rata saja?
\end{enumerate}

\begin{lstlisting}[language=bash, caption={Membaca metrik mentah lalu papannya, tersedia di \berkas{sources/chapter09/papan.py}}, label={lst:c9-lat2-jalan}]
curl -s http://127.0.0.1:8061/metrics | head -8
python3 papan.py

# Keluarannya:
#   Versi                       Minta   Gagal  % gagal   Rata ms   p95 ms
#   biru                          450       0      0.0       5.1       10
#   hijau                          50       2      4.0      20.1       25
\end{lstlisting}

\begin{table}[!htb]
  \centering
  \caption{Contoh hasil Latihan 2 yang sudah terisi}
  \label{tab:c9-lat2-contoh}
  \begin{tabularx}{\textwidth}{Xrrrr}
    \hline
    \textbf{Versi} & \textbf{Minta} & \textbf{Gagal} & \textbf{Rata ms} & \textbf{p95 ms} \\
    \hline
    biru & 450 & 0 & 5,1 & 10 \\
    \hline
    hijau & 50 & 2 & 20,1 & 25 \\
    \hline
  \end{tabularx}
\end{table}

\begin{table}[!htb]
  \centering
  \caption{Lembar isian hasil pengamatan Latihan 2}
  \label{tab:c9-lat2-hasil}
  \begin{tabularx}{\textwidth}{Xrrrr}
    \hline
    \textbf{Versi} & \textbf{Minta} & \textbf{Gagal} & \textbf{Rata ms} & \textbf{p95 ms} \\
    \hline
    biru, pengamatan pertama & \dots & \dots & \dots & \dots \\
    \hline
    hijau, pengamatan pertama & \dots & \dots & \dots & \dots \\
    \hline
    biru, sesudah 200 permintaan & \dots & \dots & \dots & \dots \\
    \hline
    hijau, sesudah 200 permintaan & \dots & \dots & \dots & \dots \\
    \hline
  \end{tabularx}
\end{table}

\subsection*{Latihan 3: Menelusuri Satu Permintaan}

Latihan ini membuktikan bahwa \textit{trace id} menjawab pertanyaan yang
tidak dapat dijawab metrik maupun \textit{log} sendirian.

\begin{enumerate}
  \item Ambil satu \textit{trace id} dari header \texttt{X-Trace-Id} sebuah
        \textit{response}, seperti pada Listing \ref{lst:c9-lat3-jalan}.
  \item Cari penanda itu di kedua file log, lalu catat status dan lamanya di
        tiap proses ke Tabel \ref{tab:c9-lat3-hasil}. Tabel
        \ref{tab:c9-lat3-contoh} memperlihatkan contoh yang sudah terisi.
  \item Ulangi untuk satu permintaan yang gagal, lalu catat di proses mana
        kegagalannya berasal.
  \item Jawab dua pertanyaan berikut. Pertama, berapa biaya penerusan
        proxy, dan dari angka mana Anda mengetahuinya? Kedua, apa yang
        terjadi pada penelusuran bila proxy tidak meneruskan header
        \texttt{traceparent}?
\end{enumerate}

\begin{lstlisting}[language=bash, caption={Mengambil trace id lalu mencarinya di kedua log, tersedia di \berkas{sources/chapter09/proxy.py}}, label={lst:c9-lat3-jalan}]
curl -sI http://127.0.0.1:8060/api/katalog | grep -i x-trace-id
grep -h "300a032b60413b55e53009cc0660f240" proxy.log hijau.log

# Keluarannya, dirapikan:
#   hijau    /api/katalog   status 500   20.2 ms
#   proxy    /api/katalog   status 500   22.2 ms
\end{lstlisting}

\begin{table}[!htb]
  \centering
  \caption{Contoh hasil Latihan 3 yang sudah terisi}
  \label{tab:c9-lat3-contoh}
  \begin{tabularx}{\textwidth}{lrr}
    \hline
    \textbf{Proses} & \textbf{Status} & \textbf{Lama} \\
    \hline
    hijau & 500 & 20,2 ms \\
    \hline
    proxy & 500 & 22,2 ms \\
    \hline
  \end{tabularx}
\end{table}

\begin{table}[!htb]
  \centering
  \caption{Lembar isian hasil penelusuran Latihan 3}
  \label{tab:c9-lat3-hasil}
  \begin{tabularx}{\textwidth}{lrr}
    \hline
    \textbf{Proses} & \textbf{Status} & \textbf{Lama} \\
    \hline
    Permintaan berhasil, di proxy & \dots & \dots \\
    \hline
    Permintaan berhasil, di layanan & \dots & \dots \\
    \hline
    Permintaan gagal, di proxy & \dots & \dots \\
    \hline
    Permintaan gagal, di layanan & \dots & \dots \\
    \hline
  \end{tabularx}
\end{table}

\subsection*{Latihan 4: Menegakkan SLO dan Menjalankan Container}

Latihan ini menutup bab dengan dua hal yang saling melengkapi: batas yang
dinilai mesin, dan artefak yang dapat dijalankan di mana saja.

\begin{enumerate}
  \item Jalankan alarmnya, lalu catat \textit{verdict} tiap tujuan beserta
        status keluarnya, seperti pada Listing \ref{lst:c9-lat4-jalan}.
  \item Turunkan \texttt{ketersediaan\_persen\_minimum} di
        \berkas{sources/chapter09/slo.json} menjadi 95, jalankan ulang
        alarmnya, lalu catat perubahannya.
  \item Bangun \textit{image} dan jalankan containernya, lalu catat ukuran
        \textit{image} serta lama sampai siap melayani ke Tabel
        \ref{tab:c9-lat4-hasil}. Tabel \ref{tab:c9-lat4-contoh}
        memperlihatkan contoh yang sudah terisi.
  \item Hapus baris \texttt{ENV HOST=0.0.0.0} dari
        \berkas{sources/chapter09/Dockerfile}, bangun ulang, lalu catat apa
        yang terjadi. Kembalikan sesudahnya.
  \item Jawab tiga pertanyaan berikut. Pertama, mengapa menurunkan SLO bukan
        cara memperbaiki layanan? Kedua, mengapa \textit{build} kedua jauh
        lebih singkat? Ketiga, mengapa container tanpa
        \texttt{HOST=0.0.0.0} tidak pernah terjawab?
\end{enumerate}

\begin{lstlisting}[language=bash, caption={Menjalankan alarm lalu membangun container, tersedia di \berkas{sources/chapter09/alarm.py}}, label={lst:c9-lat4-jalan}]
python3 alarm.py
echo $?
docker build -t awe9-layanan .
docker run -d --name awe9-uji -e VERSI=biru -p 8065:8061 awe9-layanan
curl -s http://127.0.0.1:8065/sehat

# Keluarannya:
# Gagal: 1 tujuan SLO tidak terpenuhi.
# 1
# {"versi": "biru", "keadaan": "sehat"}
\end{lstlisting}

\begin{table}[!htb]
  \centering
  \caption{Contoh hasil Latihan 4 yang sudah terisi}
  \label{tab:c9-lat4-contoh}
  \begin{tabularx}{\textwidth}{Xr}
    \hline
    \textbf{Yang diamati} & \textbf{Hasil} \\
    \hline
    Tujuan SLO yang dilanggar & 1 \\
    \hline
    Status keluar alarm & 1 \\
    \hline
    Ukuran \textit{image} & 119 MB \\
    \hline
    Lama sampai siap melayani & 0,5 s \\
    \hline
  \end{tabularx}
\end{table}

\begin{table}[!htb]
  \centering
  \caption{Lembar isian hasil pengamatan Latihan 4}
  \label{tab:c9-lat4-hasil}
  \begin{tabularx}{\textwidth}{Xr}
    \hline
    \textbf{Yang diamati} & \textbf{Hasil} \\
    \hline
    Tujuan SLO yang dilanggar & \dots \\
    \hline
    Status keluar alarm & \dots \\
    \hline
    Status keluar sesudah SLO diturunkan & \dots \\
    \hline
    Ukuran \textit{image} & \dots \\
    \hline
    Lama \textit{build} pertama dan kedua & \dots \\
    \hline
    Akibat menghapus \texttt{ENV HOST=0.0.0.0} & \dots \\
    \hline
  \end{tabularx}
\end{table}

Setelah seluruh latihan selesai, ketiga proses dihentikan dengan
\texttt{kill}, dan containernya dihapus dengan
\texttt{docker rm -f awe9-uji}.
